This section surveys the key developments in speaker recognition that inform our work, organized into three areas: the evolution of speaker embedding architectures, advances in speaker enrollment strategies, and the emerging role of phonetic information in speaker modeling.

\subsection{Speaker Embedding Architectures}

Early speaker recognition systems represented speakers using statistical models of handcrafted acoustic features. The most influential of these was the \textbf{i-vector} framework \cite{dehak2010ivector}, which employed a Gaussian Mixture Model--Universal Background Model (GMM-UBM) pipeline to compress variable-length utterances into fixed-dimensional vectors through factor analysis. Although i-vectors became a de facto standard for over a decade, their reliance on shallow generative modeling left them sensitive to channel mismatch, background noise, and speaking-style variation.

Deep neural networks have since supplanted statistical approaches. Snyder et al.\ introduced \textbf{x-vectors} \cite{snyder2018xvectors}, replacing the GMM back-end with a discriminatively trained neural architecture. The x-vector system feeds MFCC features through a stack of \textbf{Time Delay Neural Network (TDNN)} layers \cite{peddinti2015tdnn} that model temporal context at the frame level, aggregates them via a statistics pooling layer, and classifies speakers with a softmax objective. By learning representations end-to-end, x-vectors substantially outperformed i-vectors across multiple benchmarks.

Building on the TDNN backbone, Desplanques et al.\ proposed \textbf{ECAPA-TDNN} \cite{ecapa-tdnn}, which augments the x-vector paradigm with Squeeze-and-Excitation blocks for channel-wise attention, multi-scale feature aggregation, and attentive statistics pooling. Dawalatabad et al.\ \cite{dawalatabad2021ecapa} subsequently demonstrated that ECAPA-TDNN embeddings generalize well beyond verification to speaker diarization tasks. Training with the \textbf{ArcFace} angular-margin loss \cite{deng2019arcface} further sharpens inter-speaker discrimination, yielding an Equal Error Rate (EER) of 0.87\% on the \textbf{VoxCeleb1} test set \cite{nagrani2017voxceleb}. Meanwhile, Jung et al.\ pursued a complementary direction with the \textbf{RawNet} family \cite{rawnet,rawnet2}, culminating in \textbf{RawNet3} \cite{rawnet3}, which operates directly on raw waveforms and eliminates the need for hand-designed spectral features altogether. These two architectures---ECAPA-TDNN and RawNet3---represent contrasting design philosophies (engineered features vs.\ learned front-ends) and form the basis of our comparative study.

\subsection{Speaker Enrollment Strategies}

Reliable speaker identification depends not only on the embedding model but also on the quality of the enrollment data used to construct each speaker's reference profile. Several lines of work have sought to improve enrollment.

Li et al.\ \cite{li2024enrollment} proposed \textbf{enrollment speech augmentation}, applying noise and reverberation perturbations to enrollment utterances so that the resulting speaker template becomes more resilient to acoustic mismatch at test time. Mingote et al.\ \cite{mingote2020enrollment} took a model-centric approach with \textbf{network-optimized enrollment models}: instead of averaging embeddings, they trained a per-speaker output head whose weights serve as the speaker template, eliminating the need for a separate scoring back-end. Le and Ngo \cite{le2025datacentric} introduced a \textbf{data-centric enrollment selection} algorithm that iteratively retains only those utterances whose pairwise cosine similarities exceed an adaptive threshold, filtering out noisy or atypical recordings while preserving representative ones. Despite their effectiveness, all of these methods operate on the acoustic and environmental dimension of enrollment quality, largely ignoring the linguistic content of the enrolled speech.

\subsection{Phonetic Information in Speaker Modeling}

A parallel body of research has established that phonetic content plays a significant role in how well speaker characteristics are captured. Liu et al.\ \cite{liu2018phonetic,liu2019phonetic} conducted early investigations showing that \textbf{phoneme-aware speaker embeddings}---obtained by conditioning the network on phonetic labels---consistently improve verification accuracy over content-agnostic baselines. Wang et al.\ \cite{wang2019phonetic} corroborated this finding for text-independent scenarios, demonstrating that even implicit phonetic structure benefits embedding quality. Subsequent work has refined these ideas with more sophisticated architectures: Liu et al.\ \cite{liu2022phonetic} introduced a \textbf{phonetic attention mask} that guides the network to weight phonetically informative frames more heavily, while Hong et al.\ \cite{hong2023phonetic} proposed \textbf{phonetic decomposition and reorganization} modules that disentangle speaker and phonetic factors within the embedding space. Most recently, Song et al.\ \cite{song2024phonetic} extended these techniques to the \textbf{multilingual} setting, showing that cross-lingual phonetic features transfer effectively for speaker verification across languages. Complementing these neural approaches, Dellwo et al.\ \cite{dellwo2015rhythm} provided acoustic-phonetic evidence that \textbf{prosodic and rhythmic variability} across a wider phonetic range better captures individual speaker traits.

In the enrollment domain, the \textbf{TIMIT} corpus \cite{garofolo1993timit}---a phonetically balanced speech database originally designed for ASR research---illustrates the long-recognized principle that controlled phonetic coverage leads to more informative speech samples. Yet, to the best of our knowledge, no prior work has explicitly leveraged phonetically rich, guided utterances during the enrollment phase for speaker identification. Le and Ngo \cite{longquocandnhut} recently addressed this gap by proposing a \textbf{phoneme-based enrollment optimization} strategy that selects enrollment sentences to maximize phoneme coverage, thereby producing higher-quality speaker profiles. Their results demonstrate that bridging linguistic diversity with enrollment design leads to more robust and consistent recognition, and their approach provides a key foundation for the enrollment methodology adopted in this project.